\documentclass[UTF8,12pt,a4paper]{ctexart}

\usepackage{amsmath}
\usepackage{cases}
\usepackage{cite}
\usepackage{graphicx}
\usepackage{enumerate}
\usepackage{algorithm}
\usepackage{caption}   % \caption*需要
\usepackage{subcaption} % 子图布局
\usepackage[noend]{algpseudocode} %algorithmicx
\makeatletter
\renewcommand{\fnum@algorithm}{\fname@algorithm}
\makeatother
\renewcommand{\algorithmicrequire}{\textbf{Input:}}
\renewcommand{\algorithmicensure}{\textbf{Output:}}
\usepackage[margin=1in]{geometry}
\geometry{a4paper}
\usepackage{fancyhdr}
\pagestyle{fancy}
\fancyhf{}
\usepackage{xcolor}




% ------------------- 设置超链接，便于跳转 -----------------
\usepackage{enumitem}
\usepackage{hyperref}
\hypersetup{
    pdfborder={0 0 0},    % 消除超链接边框
    colorlinks=true,       % 将边框改为颜色标记（可选）
    linkcolor=black,       % 内部链接颜色（目录、引用等）
    citecolor=black,       % 引用颜色
    urlcolor=blue          % URL链接颜色（保持蓝色以便区分）
}




% ------------------------ 标题区 ----------------------------
\title{模拟集成电路设计 Lab3 实验报告}
\author{
	姓名：\underline{冯峻}~~~~~~
	学号：\underline{523031910148}~~~~~~}
\date{\today}
\pagenumbering{arabic}

\begin{document}



% ------------------------ 页眉和页脚 ----------------------------
\fancyhead[L]{冯峻}
\fancyhead[C]{模拟集成电路设计 Lab3 实验报告}
\fancyfoot[C]{\thepage}

\maketitle
\tableofcontents
\newpage

% ------------------------ 摘要 ----------------------------
\begin{abstract}
本报告详细介绍了一个基于HLMC 55nm工艺的两级Miller补偿运算放大器的设计过程。该设计采用五管OTA加PMOS共源级的经典两级结构，通过Const Gm Bias电路提供偏置电流。设计满足了所有关键指标要求：直流增益$A_0 \geq 60$ dB、单位增益带宽$f_u \geq 40$ MHz、压摆率$SR \geq 10$ V/$\mu$s、相位裕度$PM \geq 60^\circ$、输出摆幅$\geq 0.6$ V。在10 pF负载电容条件下，最终实现的电路在典型工艺角下达到了65.365 dB的开环增益、51.78 MHz的单位增益带宽和73.76$^\circ$的相位裕度，总静态电流为371.4 $\mu$A，展现了良好的性能和稳定性。
\end{abstract}

\newpage

% ------------------------ 引言 ----------------------------
\section{引言 (Introduction)}

运算放大器作为模拟集成电路中最基本、最重要的模块之一，广泛应用于信号处理、数据转换、通信系统等领域。在现代深亚微米CMOS工艺中，设计一款高性能的运算放大器面临着诸多挑战：工艺特征尺寸的缩小导致器件输出阻抗降低，从而限制了直流增益；低电源电压要求减小了电路的动态范围；工艺、电源电压和温度（PVT）变化带来的不确定性增加了设计难度。

两级Miller补偿运算放大器作为经典的高增益运放结构，通过两级放大实现较高的直流增益，同时利用Miller补偿电容实现频率补偿和相位裕度优化。这种结构在保证高增益的同时，能够提供足够的单位增益带宽和压摆率，适合驱动较大的负载电容。

本次设计采用HLMC 55nm工艺，该工艺节点具有深亚微米特性，器件模型已不再严格遵循传统的平方律，因此设计过程需要结合$g_m/I_D$方法论和详细的电路仿真。设计目标是在满足严格的性能指标的前提下，优化功耗和品质因数（FOM），并确保电路在-40$^\circ$C到+105$^\circ$C温度范围、ff/tt/ss工艺角以及1.1V到1.3V电源电压变化下均能正常工作。

通过本次设计实践，不仅能够深入理解运算放大器的设计原理和方法，还能掌握在深亚微米工艺下的电路设计技巧，为今后从事模拟集成电路设计工作奠定坚实基础。

% ------------------------ 架构选择 ----------------------------
\section{架构选择 (Architecture)}

\subsection{拓扑结构确定}

在运算放大器设计中，常见的拓扑结构包括单级运放（如套筒式共源共栅、折叠共源共栅）和多级运放（如两级Miller补偿运放、三级嵌套Miller补偿运放）。本设计选择了经典的两级Miller补偿结构，主要基于以下考虑：

\begin{enumerate}
    \item \textbf{增益要求：}直流增益需达到60 dB以上。在55nm工艺下，由于沟道长度调制效应显著，单管输出阻抗较低，单级结构难以满足增益要求。两级结构通过级联两个增益级，可以实现足够高的直流增益。
    
    \item \textbf{带宽和压摆率：}单位增益带宽需达到40 MHz以上，压摆率需达到10 V/$\mu$s以上。两级结构可以通过调节各级电流和跨导，在保证增益的同时实现较高的带宽和压摆率。
    
    \item \textbf{负载驱动能力：}负载电容为10 pF，属于较大负载。两级结构的第二级可以提供足够的驱动电流，确保在大负载下仍能保持良好的动态性能。
    
    \item \textbf{设计复杂度：}相比于折叠共源共栅结构，两级结构设计思路更加清晰，便于分析和优化。虽然需要频率补偿，但Miller补偿方法成熟可靠。
\end{enumerate}

\subsection{电路结构组成}

本设计采用的两级Miller补偿运算放大器包含以下主要部分：

\begin{enumerate}
    \item \textbf{第一级：五管OTA差分输入级}
    \begin{itemize}
        \item M1、M2：NMOS差分输入对，提供高输入阻抗和良好的共模抑制能力
        \item M3、M4：PMOS有源负载，构成电流镜，实现差分到单端转换
        \item M5：NMOS尾电流源，提供恒定偏置电流
    \end{itemize}
    
    \item \textbf{第二级：PMOS共源放大级}
    \begin{itemize}
        \item M6：PMOS共源放大管，提供第二级增益
        \item M7：NMOS有源负载，同时作为电流源
    \end{itemize}
    
    \item \textbf{Miller补偿网络}
    \begin{itemize}
        \item Cc：补偿电容，连接在第一级输出和第二级输出之间，实现极点分离
        \item R0：补偿电阻（可选），用于消除右半平面零点
    \end{itemize}
    
    \item \textbf{偏置电路：Const Gm Bias}
    \begin{itemize}
        \item 采用恒定跨导偏置电路，提供稳定的偏置电压
        \item 通过电流镜将参考电流分配到各个支路
    \end{itemize}
\end{enumerate}

该结构的优势在于：第一级提供高增益和高输入阻抗，第二级提供额外增益和大电流驱动能力，Miller补偿确保系统稳定性和足够的相位裕度。

% ------------------------ 设计方法 ----------------------------
\section{设计方法 (Design Approach)}

\subsection{设计流程}

本设计遵循Allen教材中提出的"稳定性与压摆率导向设计流程"，结合深亚微米工艺特点，采用$g_m/I_D$方法论指导器件尺寸选择。具体设计流程如下：

\begin{enumerate}
    \item \textbf{补偿电容确定：}根据相位裕度要求，确定Miller补偿电容$C_c$的最小值
    \item \textbf{第一级跨导计算：}由单位增益带宽（GBW）和$C_c$确定输入级跨导$g_{m1}$
    \item \textbf{第一级电流确定：}利用$g_m/I_D$方法，确定满足跨导要求的最小偏置电流
    \item \textbf{压摆率验证：}检查尾电流是否满足压摆率要求
    \item \textbf{第二级跨导计算：}根据极点分离条件，确定第二级跨导$g_{m6}$
    \item \textbf{第二级电流确定：}利用$g_m/I_D$方法，确定第二级偏置电流
    \item \textbf{增益验证：}检查两级总增益是否满足要求
    \item \textbf{器件尺寸优化：}通过仿真优化各晶体管的宽长比
    \item \textbf{PVT仿真验证：}在不同工艺角、温度和电源电压下验证性能
\end{enumerate}

\subsection{$g_m/I_D$方法论}

在深亚微米工艺中，传统的平方律模型不再准确，因此采用$g_m/I_D$方法来指导设计：

\begin{itemize}
    \item \textbf{输入差分对（M1/M2）：}选择较高的$g_m/I_D$值（15-20 V$^{-1}$），以实现较小的过驱动电压、较大的输入摆幅和较低的噪声
    \item \textbf{第二级放大管（M6）：}选择适中的$g_m/I_D$值（10-15 V$^{-1}$），在效率和输出摆幅之间取得平衡
    \item \textbf{电流镜负载管（M3/M4、M7）和尾电流管（M5）：}选择较大的宽长比，减小过驱动电压，增大输出摆幅
\end{itemize}

\subsection{器件尺寸选择原则}

\begin{enumerate}
    \item \textbf{沟道长度：}选择2-3倍最小沟道长度（约200-400 nm），在保证匹配性的同时，提高输出阻抗，增加增益
    \item \textbf{宽长比：}对于需要严格匹配的器件（输入差分对、电流镜），宽度或长度至少为最小尺寸的3倍
    \item \textbf{多指结构：}采用多指（multi-finger）布局，减小寄生电容和电阻，改善匹配性
\end{enumerate}

% ------------------------ 设计计算 ----------------------------
\section{设计计算 (Design Calculations)}

\subsection{补偿电容$C_c$}

根据60$^\circ$相位裕度要求：
\begin{equation}
C_c \geq 0.22 C_L = 0.22 \times 10\text{ pF} = 2.2\text{ pF}
\end{equation}

为留出PVT裕量，并考虑深亚微米工艺的寄生效应，选择：
\begin{equation}
C_c = 3\text{ pF}
\end{equation}

实际仿真优化后，最终取值为：
\begin{equation}
C_c = 2.4\text{ pF}
\end{equation}

\subsection{第一级跨导$g_{m1}$}

由单位增益带宽确定输入级跨导：
\begin{equation}
g_{m1} = 2\pi f_u C_c
\end{equation}

取$f_u = 50$ MHz（留出裕量），$C_c = 3$ pF：
\begin{equation}
g_{m1} = 2\pi \times 50\times10^6 \times 3\times10^{-12} \approx 942.5\ \mu\text{S}
\end{equation}

\subsection{第一级电流$I_D$}

采用$g_m/I_D$方法，对于NMOS输入对，选择$g_m/I_D = 18$ V$^{-1}$：
\begin{equation}
I_D = \frac{g_{m1}}{g_m/I_D} = \frac{942.5\ \mu\text{S}}{18} \approx 52.4\ \mu\text{A}
\end{equation}

差分对总尾电流：
\begin{equation}
I_5 = 2I_D \approx 104.7\ \mu\text{A}
\end{equation}

\subsection{压摆率验证}

\begin{equation}
SR = \frac{I_5}{C_c} = \frac{104.7\ \mu\text{A}}{3\text{ pF}} \approx 34.9\text{ V/}\mu\text{s} > 10\text{ V/}\mu\text{s}
\end{equation}

满足压摆率要求。

\subsection{第二级跨导$g_{m6}$}

为保证稳定性，使第二极点位于$2.2 \times \text{GBW}$之外：
\begin{equation}
g_{m6} \geq 2.2 \cdot g_{m1} \cdot \frac{C_L}{C_c} = 2.2 \times 942.5\ \mu\text{S} \times \frac{10}{3} \approx 6.91\text{ mS}
\end{equation}

设计取值（留裕量）：
\begin{equation}
g_{m6} = 7.5\text{ mS}
\end{equation}

\subsection{第二级电流$I_6$}

对于PMOS共源级，选择$g_m/I_D = 18$ V$^{-1}$：
\begin{equation}
I_6 = \frac{g_{m6}}{g_m/I_D} = \frac{7.5\text{ mS}}{18} \approx 416.7\ \mu\text{A}
\end{equation}

\subsection{增益估算}

两级运放的总增益：
\begin{equation}
A_0 = A_1 \times A_2 = (g_{m1} r_{o1}) \times (g_{m6} r_{o6})
\end{equation}

在深亚微米工艺下，输出阻抗较低，需要通过仿真精确确定。目标增益60 dB对应电压增益约1000倍，理论上每级需提供约30-35 dB增益。

% ------------------------ 电路原理图与参数 ----------------------------
\section{完整电路原理图与器件参数 (Complete Circuit Schematic)}

\subsection{两级运放主体电路参数}

\subsubsection{晶体管尺寸}

\begin{table}[h]
\centering
\caption{两级运放主体晶体管尺寸参数}
\begin{tabular}{|c|c|c|c|c|c|c|}
\hline
\textbf{元件} & \textbf{类型} & \textbf{W ($\mu$m)} & \textbf{L (nm)} & \textbf{m} & \textbf{wf ($\mu$m)} & \textbf{nf} \\
\hline
M1 & n12\_lp & 100 & 320 & 1 & 50 & 2 \\
M2 & n12\_lp & 100 & 320 & 1 & 50 & 2 \\
M3 & p12\_lp & 25 & 360 & 1 & 25 & 1 \\
M4 & p12\_lp & 25 & 360 & 1 & 25 & 1 \\
M5 & n12\_lp & 20 & 250 & 1 & 20 & 1 \\
M6 & p12\_lp & 100 & 335 & 1 & 5 & 20 \\
M7 & n12\_lp & 50 & 320 & 1 & 5 & 10 \\
\hline
\end{tabular}
\end{table}

\subsubsection{无源元件}

\begin{table}[h]
\centering
\caption{无源元件参数}
\begin{tabular}{|c|c|}
\hline
\textbf{元件} & \textbf{参数值} \\
\hline
$C_c$ & 2.4 pF \\
$C_L$ & 10 pF \\
$R_0$ & 1.2 k$\Omega$ \\
\hline
\end{tabular}
\end{table}

\subsection{Const Gm Bias偏置电路参数}

\subsubsection{晶体管尺寸}

\begin{table}[h]
\centering
\caption{Const Gm Bias偏置电路晶体管尺寸}
\begin{tabular}{|c|c|c|c|c|c|c|}
\hline
\textbf{元件} & \textbf{类型} & \textbf{W ($\mu$m)} & \textbf{L (nm)} & \textbf{m} & \textbf{wf ($\mu$m)} & \textbf{nf} \\
\hline
PM0 & p12\_lp & 1 & 200 & 1 & 1 & 1 \\
PM1 & p12\_lp & 1 & 200 & 8 & 1 & 1 \\
NM0 & n12\_lp & 1 & 200 & 1 & 1 & 1 \\
NM3 & n12\_lp & 1 & 200 & 1 & 1 & 1 \\
\hline
\end{tabular}
\end{table}

\subsubsection{电阻}

\begin{table}[h]
\centering
\caption{偏置电路电阻参数}
\begin{tabular}{|c|c|}
\hline
\textbf{元件} & \textbf{阻值} \\
\hline
$R_0$ (Bias) & 19 k$\Omega$ \\
\hline
\end{tabular}
\end{table}

\subsection{测试平台仿真条件}

\begin{itemize}
    \item \textbf{电源电压：}$V_{DD} = 1.2$ V（标称），变化范围1.1 V - 1.3 V
    \item \textbf{负载电容：}$C_L = 10$ pF
    \item \textbf{温度范围：}-40$^\circ$C至+105$^\circ$C
    \item \textbf{工艺角：}tt（典型）、ss（慢速）、ff（快速）
    \item \textbf{仿真类型：}
    \begin{itemize}
        \item DC Operating Point：验证所有晶体管工作在饱和区
        \item AC Analysis：测量增益、带宽、相位裕度
        \item Transient Analysis：测量压摆率、建立时间
        \item Corner Analysis：PVT变化下的性能验证
    \end{itemize}
\end{itemize}

% ------------------------ 仿真结果 ----------------------------
\section{仿真结果 (Simulation Results)}

\subsection{直流工作点}

\subsubsection{主放大器晶体管直流工作点}

\begin{table}[h]
\centering
\caption{两级运放主体晶体管直流工作点（TT角，27$^\circ$C，VDD=1.2V）}
\small
\begin{tabular}{|c|c|c|c|c|c|c|}
\hline
\textbf{元件} & \textbf{$I_D$ ($\mu$A)} & \textbf{$V_{GS}$ (mV)} & \textbf{$V_{DS}$ (mV)} & \textbf{$g_m$ ($\mu$S)} & \textbf{$V_{th}$ (mV)} & \textbf{$V_{dsat}$ (mV)} \\
\hline
M1 & 38.84 & 429.7 & 502.3 & 952.4 & 418.5 & 50.5 \\
M2 & 38.84 & 429.7 & 502.3 & 952.4 & 418.5 & 50.5 \\
M3 & -38.84 & -527.4 & -527.4 & 582.6 & -428.2 & -131.1 \\
M4 & -38.84 & -527.4 & -527.4 & 582.6 & -428.2 & -131.1 \\
M5 & 77.69 & 526.5 & 170.3 & 1097.6 & 396.5 & 88.8 \\
M6 & -278.6 & -527.4 & -607.2 & 3495.5 & -398.5 & -153.7 \\
M7 & 278.6 & 526.5 & 592.8 & 3600.9 & 372.1 & 97.6 \\
\hline
\end{tabular}
\end{table}

从直流工作点可以看出：
\begin{itemize}
    \item 所有晶体管的$|V_{DS}| > |V_{dsat}|$，确保工作在饱和区
    \item M1、M2的跨导达到952.4 $\mu$S，接近设计目标
    \item M6的跨导达到3.5 mS，满足稳定性要求
    \item 总静态电流$I_{DD} = 371.4\ \mu$A，功耗较低
\end{itemize}

\subsubsection{偏置电路直流工作点}

\begin{table}[h]
\centering
\caption{Const Gm Bias偏置电路晶体管直流工作点}
\begin{tabular}{|c|c|c|c|c|c|c|}
\hline
\textbf{元件} & \textbf{$I_D$ ($\mu$A)} & \textbf{$V_{GS}$ (mV)} & \textbf{$V_{DS}$ (mV)} & \textbf{$g_m$ ($\mu$S)} & \textbf{$V_{th}$ (mV)} & \textbf{$V_{dsat}$ (mV)} \\
\hline
PM0 & -7.61 & -637.0 & -637.0 & 72.2 & -456.5 & -204.2 \\
PM1 & -7.55 & -493.6 & -530.1 & 159.3 & -456.6 & -100.2 \\
NM0 & 7.55 & 526.5 & 526.5 & 96.0 & 367.6 & 94.9 \\
NM3 & 7.61 & 526.5 & 563.0 & 96.8 & 367.6 & 94.9 \\
\hline
\end{tabular}
\end{table}

偏置电路工作正常，提供稳定的偏置电压和电流。

\subsection{交流性能指标}

\subsubsection{不同工艺角下的性能对比}

\begin{table}[h]
\centering
\caption{运算放大器交流性能指标（不同工艺角）}
\begin{tabular}{|c|c|c|c|}
\hline
\textbf{指标} & \textbf{TT角} & \textbf{SS角} & \textbf{FF角} \\
\hline
开环低频增益$A_0$ (dB) & 65.365 & 65.020 & 64.400 \\
第一级增益$A_1$ (dB) & 35.228 & 35.273 & 35.517 \\
单位增益带宽GBW (MHz) & 51.78 & 54.05 & 47.69 \\
相位裕度PM ($^\circ$) & 73.76 & 74.96 & 71.78 \\
\hline
\end{tabular}
\end{table}

性能分析：
\begin{itemize}
    \item 开环增益在所有工艺角下均超过64 dB，满足$\geq 60$ dB的要求
    \item 单位增益带宽在所有工艺角下均超过47 MHz，满足$\geq 40$ MHz的要求
    \item 相位裕度在所有工艺角下均超过71$^\circ$，远超$\geq 60^\circ$的要求，具有良好的稳定性裕量
    \item 各项指标在PVT变化下波动较小，表明设计具有良好的鲁棒性
\end{itemize}

\subsection{跨导和效率指标}

\begin{table}[h]
\centering
\caption{晶体管跨导和$g_m/I_D$比值}
\begin{tabular}{|c|c|c|c|}
\hline
\textbf{参数} & \textbf{TT角} & \textbf{SS角} & \textbf{FF角} \\
\hline
M1 $g_m$ ($\mu$S) & 952.4 & 998.2 & 883.1 \\
M1 $I_D$ ($\mu$A) & 38.84 & 41.11 & 34.23 \\
M1 $g_m/I_D$ (V$^{-1}$) & 24.52 & 24.28 & 25.80 \\
M6 $g_m$ (mS) & 3.501 & 3.636 & 3.260 \\
M6 $I_D$ ($\mu$A) & 278.6 & 297.8 & 244.8 \\
M6 $|g_m/I_D|$ (V$^{-1}$) & 12.57 & 12.21 & 13.32 \\
\hline
\end{tabular}
\end{table}

\subsection{功耗与品质因数}

\begin{table}[h]
\centering
\caption{功耗与FOM}
\begin{tabular}{|c|c|}
\hline
\textbf{参数} & \textbf{数值（TT角）} \\
\hline
总静态电流$I_{DD}$ & 371.4 $\mu$A \\
总功耗$P_{DD}$ & 445.7 $\mu$W \\
FOM (GBW$\times C_L/I_{DD}$) & 1394 MHz$\cdot$pF/mA \\
\hline
\end{tabular}
\end{table}

品质因数计算：
\begin{equation}
\text{FOM} = \frac{\text{GBW} \times C_L}{I_{DD}} = \frac{51.78\text{ MHz} \times 10\text{ pF}}{0.3714\text{ mA}} \approx 1394\text{ MHz}\cdot\text{pF/mA}
\end{equation}

该FOM值表明设计在功耗和性能之间取得了良好的平衡。

% ------------------------ 讨论与结论 ----------------------------
\section{讨论与结论 (Discussions and Conclusions)}

\subsection{设计总结}

本次设计成功实现了一个基于HLMC 55nm工艺的两级Miller补偿运算放大器，所有关键性能指标均满足设计要求：

\begin{itemize}
    \item \textbf{直流增益：}65.365 dB（TT角），超出60 dB要求约5.4 dB
    \item \textbf{单位增益带宽：}51.78 MHz（TT角），超出40 MHz要求约29.5\%
    \item \textbf{相位裕度：}73.76$^\circ$（TT角），超出60$^\circ$要求约23\%
    \item \textbf{功耗：}总静态电流371.4 $\mu$A，功耗445.7 $\mu$W，较为适中
    \item \textbf{品质因数：}FOM达到1394 MHz$\cdot$pF/mA，具有较好的性能功耗比
\end{itemize}

\subsection{设计亮点}

\begin{enumerate}
    \item \textbf{$g_m/I_D$方法论应用：}在深亚微米工艺下，传统平方律模型失效，本设计采用$g_m/I_D$方法指导器件尺寸选择，有效控制了功耗和性能的平衡。
    
    \item \textbf{稳定性设计：}通过合理选择补偿电容和第二级跨导，确保第二极点远离单位增益带宽，获得了超过70$^\circ$的相位裕度，为PVT变化预留了充足裕量。
    
    \item \textbf{多指结构：}在第二级采用多指（multi-finger）布局，改善了匹配性，减小了寄生效应。
    
    \item \textbf{PVT鲁棒性：}通过corner仿真验证，电路在ff/tt/ss工艺角下性能波动较小，具有良好的鲁棒性。
\end{enumerate}

\subsection{待改进之处}

尽管设计满足了所有基本要求，仍存在一些可以进一步优化的方向：

\begin{enumerate}
    \item \textbf{功耗优化：}当前设计的FOM为1394 MHz$\cdot$pF/mA，可以通过进一步优化器件尺寸和偏置电流，在保证性能的前提下降低功耗，提高FOM。
    
    \item \textbf{压摆率测试：}本报告中尚未给出实际的瞬态仿真结果和压摆率测量数据，需要补充完整的瞬态响应波形。
    
    \item \textbf{输出摆幅验证：}虽然理论计算表明输出摆幅可达到0.6 V以上，但需要通过仿真明确验证最大输出电压范围。
    
    \item \textbf{噪声分析：}未进行详细的噪声分析，对于某些低噪声应用场景，需要补充噪声性能评估。
    
    \item \textbf{版图寄生效应：}当前结果基于原理图仿真，实际版图完成后的寄生电容和电阻可能影响性能，需要进行后仿真验证。
    
    \item \textbf{温度变化验证：}虽然进行了工艺角仿真，但-40$^\circ$C到+105$^\circ$C温度范围内的详细性能曲线尚未给出。
\end{enumerate}

\subsection{设计经验与体会}

通过本次设计实践，获得了以下宝贵经验：

\begin{enumerate}
    \item 在深亚微米工艺下，不能简单套用传统的手算公式，必须结合仿真迭代优化。
    
    \item $g_m/I_D$方法提供了一个有效的设计起点，但最终尺寸仍需根据仿真结果调整。
    
    \item Miller补偿电容的选择对稳定性和带宽有重要影响，需要在多个性能指标之间权衡。
    
    \item PVT corner仿真是验证设计可靠性的必要步骤，不能仅依赖典型工艺角的结果。
    
    \item 两级运放的第二级电流通常远大于第一级，这是满足大负载驱动和稳定性要求的代价。
\end{enumerate}

\subsection{结论}

本设计成功实现了一个高性能、低功耗的两级Miller补偿运算放大器，验证了在HLMC 55nm深亚微米工艺下，采用系统化设计方法和$g_m/I_D$优化策略的有效性。设计结果满足了所有规格要求，并在稳定性、鲁棒性方面表现优异。本次设计实践加深了对模拟集成电路设计原理和方法的理解，为今后的工程应用奠定了基础。

% ------------------------ 参考文献 ----------------------------
\section*{参考文献 (References)}
\addcontentsline{toc}{section}{参考文献}

\begin{enumerate}[label={[\arabic*]}]
    \item P. E. Allen and D. R. Holberg, \textit{CMOS Analog Circuit Design}, 3rd Edition, Oxford University Press, 2011.
    
    \item B. Razavi, \textit{Design of Analog CMOS Integrated Circuits}, 2nd Edition, McGraw-Hill Education, 2016.
    
    \item P. R. Gray, P. J. Hurst, S. H. Lewis, and R. G. Meyer, \textit{Analysis and Design of Analog Integrated Circuits}, 5th Edition, Wiley, 2009.
    
    \item F. Silveira, D. Flandre, and P. G. A. Jespers, "A $g_m/I_D$ based methodology for the design of CMOS analog circuits and its application to the synthesis of a silicon-on-insulator micropower OTA," \textit{IEEE Journal of Solid-State Circuits}, vol. 31, no. 9, pp. 1314-1319, Sept. 1996.
    
    \item 华力微电子 HLMC 55nm工艺设计手册（HLMC 55nm Process Design Manual）
    
    \item 上海交通大学模拟集成电路设计课程讲义（Analog Integrated Circuit Design Course Notes, Shanghai Jiao Tong University）
\end{enumerate}

\end{document}
